\documentclass[../../../../main.tex]{subfiles}
\begin{document}

\begin{flushright}
\footnotesize\textit{【自動文字起こし・要確認】}
\end{flushright}

\noindent
[\textbf{解}]
\[
\begin{aligned}
|\vec{X}|^2 &= |\vec{A}|^2 + |\vec{B_k}|^2 + 2\vec{A} \cdot \vec{B_k} \\
&= a^2 + 1 + 2\vec{A} \cdot \vec{B_k} \quad \dots \text{①'}
\end{aligned}
\]
である。題意から，$\vec{A} = a \begin{pmatrix} \cos\theta \\ \sin\theta \end{pmatrix}$ とおける。($0 \leqq \theta < 2\pi$)

\vspace{0.5em}
\noindent
$|\vec{X}|^2$ の大きさの期待値を $E$ とおくと①'から，
\[
\begin{aligned}
E &= \frac{1}{6} \sum_{k=1}^6 \left\{ a^2+1 + 2 \begin{pmatrix} \cos\theta \\ \sin\theta \end{pmatrix} \cdot \begin{pmatrix} \cos\frac{k\pi}{3} \\ \sin\frac{k\pi}{3} \end{pmatrix} \right\} \\
&= (a^2+1) + \frac{1}{3} \sum_{k=1}^6 \cos\left(\frac{k\pi}{3} - \theta\right) \quad \dots \text{②}
\end{aligned}
\]
ここで，単位円に内接する正六角形を考えると，中心から各頂点へのベクトルの $x$ 成分は $\cos\left(\frac{k\pi}{3} - \theta\right)$ ($k=1, 2, \dots, 6$) で表わされ，又，この和は零だから，
\[
\sum_{k=1}^6 \cos\left(\frac{k\pi}{3} - \theta\right) = 0
\]
②に代入して
\[
E = a^2 + 1
\]

\newpage

\end{document}
